\banner{From light to yield: the quantitative chain of the crop}\label{ch:map2-deep}
\section{Thermal time in practice}\label{sec:dm-gdd}
The development stage of Equation~\ref{eq:dm-tsum} advances with the \kw{thermal time}, the sum of the daily mean temperatures above a base temperature $T_b$. Below $T_b$ the crop does not develop; above an upper limit (about 30 $^\circ$C for temperate cereals) the rate stops rising. The stage is the ratio of the accumulated sum to the sum required for the phase:
\begin{equation}\label{eq:dm-ds}\boxed{DS=\frac{TS}{TS_{\mathrm{req}}},\qquad TS=\sum_{d}\min\bigl(\max(T_d-T_b,0),\,T_{u}-T_b\bigr)}\end{equation}
Daisy uses $DS=0$ at emergence, $DS=1$ at flowering and $DS=2$ at maturity, each phase with its own requirement. Because the sum counts only degrees above $T_b$, a cool spring delays the crop more than the fall in mean temperature suggests.

\exbox{A week of thermal time}{The daily mean temperatures of a spring week are 12, 14, 17, 9, $-1$, 6 and 11 $^\circ$C. Find the thermal time accumulated by a maize crop with $T_b=4$ $^\circ$C. How many such weeks does it need for a vegetative phase of 700 $^\circ$C\,d?}\sol
\givenline{$T_d=12,14,17,9,-1,6,11$ $^\circ$C ; $T_b=4$ $^\circ$C ; $TS_{\mathrm{req}}=700$ $^\circ$C\,d}
\wkt\ each day contributes $\max(T_d-T_b,\,0)$ to the thermal time (Equation~\ref{eq:dm-ds}). The contributions are
\[(12-4)+(14-4)+(17-4)+(9-4)+0+(6-4)+(11-4)=8+10+13+5+0+2+7\]
\[\therefore\quad TS=45\ ^\circ\text{C\,d per week}\Ans\]
Note that the day at $-1$ $^\circ$C contributes nothing, not a negative amount.
\sh{Number of such weeks for the vegetative phase}
\[\text{Weeks}=\frac{TS_{\mathrm{req}}}{TS}=\frac{700}{45}=15.6\Ans\]
\sh{Wheat at a steady 15 $^\circ$C}
With $T_b=0$ $^\circ$C each day contributes 15 $^\circ$C\,d, so flowering (1100 $^\circ$C\,d) comes after
\[\frac{1100}{15}=73\ \text{days}\Ans\]

\section{Light interception and radiation-use efficiency}\label{sec:dm-rue}
The light that reaches a depth in the canopy where the cumulative leaf area above is $L$ follows the exponential law of Monsi and Saeki (Equation~\ref{eq:dm-beer}). The fraction intercepted by the whole canopy, $f=1-e^{-kL}$, rises quickly at first and then saturates: the last leaves add little. A useful summary of production is the \kw{radiation-use efficiency} (RUE), the dry matter formed per unit of intercepted photosynthetically active radiation (PAR):
\begin{equation}\label{eq:dm-rueq}\boxed{\frac{\mathrm dW}{\mathrm dt}=\mathrm{RUE}\times f\times\mathrm{PAR},\qquad f=1-e^{-kL}}\end{equation}
For C$_3$ cereals RUE is about 2.5 to 3.0 g\,MJ$^{-1}$ of intercepted PAR, and for maize (C$_4$) about 3.5 to 4.0 g\,MJ$^{-1}$. Daisy does not use RUE; it computes the photosynthesis leaf layer by leaf layer, but the result can be checked against Equation~\ref{eq:dm-rueq}.

\begin{kfig}{Light in a canopy of four leaf layers: the flux falls exponentially with the cumulative leaf area $L$ above, $I=I_0e^{-kL}$; the fraction intercepted by the canopy is $1-e^{-kL_{\mathrm{tot}}}$.}
\foreach \y/\l in {3.2/1,2.4/2,1.6/3,0.8/4}{\draw[draw=green!45!black,line width=1.2pt] (0.3,\y)--(3.2,\y);\node[lab,left] at (0.25,\y) {$L=\l$};}
\foreach \x in {0.7,1.4,2.1,2.8}{\draw[vec] (\x,4.2)--(\x,3.35);}
\foreach \x in {0.9,1.9,2.6}{\draw[vec] (\x,3.1)--(\x,2.55);}\foreach \x in {1.4,2.4}{\draw[vec] (\x,2.3)--(\x,1.75);}\draw[vec] (1.9,1.5)--(1.9,0.95);
\draw[soil] (0,0) rectangle (3.5,0.4);
\begin{scope}[xshift=5cm]
\draw[axis] (0,0)--(3.6,0);\draw[axis] (0,0)--(0,4.4);\node[lab,below] at (1.8,-0.1) {$I/I_0$};\node[lab,rotate=90] at (-0.35,2.1) {cumulative $L$ (down)};
\draw[crv,smooth,domain=0:4] plot ({3.3*exp(-0.6*\x)},{4-\x});
\node[lab] at (2.6,3.3) {$e^{-kL}$};
\end{scope}
\end{kfig}

\exbox{Interception and growth rate of a wheat canopy}{A wheat canopy has $k=0.6$ and $L=3$. Find the fraction of light intercepted and the leaf area needed to intercept 95 per cent. With RUE $=3.0$ g\,MJ$^{-1}$ and PAR $=10$ MJ\,m$^{-2}$\,d$^{-1}$, find the daily growth.}\sol
\givenline{$k=0.6$ ; $L=3$ ; RUE $=3.0$ g\,MJ$^{-1}$ ; PAR $=10$ MJ\,m$^{-2}$\,d$^{-1}$}
\sh{1. Fraction intercepted}
\wkt\ the fraction of light intercepted by a canopy is
\[f=1-e^{-kL}\ti{i}\]
\[\therefore\quad f=1-e^{-0.6\times3}=1-e^{-1.8}=1-0.165=0.835\Ans\]
\sh{2. Leaf area for 95 per cent interception}
Putting $f=0.95$ in equation \rfi{i},
\[e^{-0.6L}=1-0.95=0.05\qquad\text{or}\qquad L=\frac{\ln(1/0.05)}{0.6}=\frac{\ln20}{0.6}=\frac{2.996}{0.6}=5.0\Ans\]
Raising the interception from 83.5 to 95 per cent therefore needs two more units of leaf area: the last leaves add little light.
\sh{3. Daily growth}
From Equation~\ref{eq:dm-rueq},
\[\frac{\mathrm dW}{\mathrm dt}=\text{RUE}\times f\times\text{PAR}=3.0\times0.835\times10=25.0\ \text{g\,m}^{-2}\,\text{d}^{-1}\Ans\]
\[=25.0\times10=250\ \text{kg DM\,ha}^{-1}\,\text{d}^{-1}\tb{$\because$ 1 g\,m$^{-2}$ = 10 kg\,ha$^{-1}$}\]

\section{The light-response curve of a leaf}\label{sec:dm-lrc}
The gross photosynthesis of a leaf rises linearly with light at low irradiance, with the \kw{initial light-use efficiency} $\varepsilon$, and saturates at $A_{\max}$ (Equation~\ref{eq:dm-light}). The curve explains two facts of the canopy. Sunlit leaves at the top work near saturation and waste part of their light; shaded leaves below work on the linear part and use every photon well. A canopy therefore produces more when the same light is spread over more leaf area, which is why its response to light is nearly linear even when each leaf saturates.

\begin{kfig}{Light-response curve of a leaf, $A=A_{\max}(1-e^{-\varepsilon I/A_{\max}})$: initial slope $\varepsilon$, saturation at $A_{\max}$; the shaded leaf (open circle) works on the linear part, the sunlit leaf (filled) near saturation.}
\draw[axis] (0,0)--(6.6,0);\draw[axis] (0,0)--(0,4.0);
\node[lab,below] at (3.3,-0.1) {absorbed light $I$ (W\,m$^{-2}$)};\node[lab,rotate=90] at (-0.45,2) {gross photosynthesis $A$};
\draw[crv,smooth,domain=0:6.2] plot (\x,{3.3*(1-exp(-1.1*\x/3.3*1.8))});
\draw[gd] (0,3.3)--(6.2,3.3);\node[lab,left] at (0,3.3) {$A_{\max}$};
\draw[gd] (0,0)--(1.7,3.2);\node[lab] at (1.25,2.6) {$\varepsilon$};
\draw[khmag,fill=white] (0.55,{3.3*(1-exp(-1.1*0.55/3.3*1.8))}) circle (0.09);\fill[khmag] (4.6,{3.3*(1-exp(-1.1*4.6/3.3*1.8))}) circle (0.09);
\node[lab,right] at (0.65,0.9) {shade leaf};\node[lab,below] at (4.6,3.15) {sunlit leaf};
\end{kfig}

\exbox{Sunlit and shaded leaves}{A leaf has $A_{\max}=40$ kg CO$_2$\,ha$^{-1}$\,h$^{-1}$ and $\varepsilon=0.45$ kg CO$_2$\,ha$^{-1}$\,h$^{-1}$ per W\,m$^{-2}$. Find its photosynthesis at 200 W\,m$^{-2}$ (sunlit) and at 50 W\,m$^{-2}$ (shaded). Which leaf uses its light more efficiently?}\sol
\givenline{$A_{\max}=40$ kg CO$_2$\,ha$^{-1}$\,h$^{-1}$ ; $\varepsilon=0.45$ kg CO$_2$\,ha$^{-1}$\,h$^{-1}$ per W\,m$^{-2}$ ; $I=200$ and $50$ W\,m$^{-2}$}
\wkt\ the light-response curve of a leaf (Equation~\ref{eq:dm-light}) is
\[A=A_{\max}\bigl(1-e^{-\varepsilon I/A_{\max}}\bigr)\ti{i}\]
\sh{1. Sunlit leaf}
\[\frac{\varepsilon I}{A_{\max}}=\frac{0.45\times200}{40}=2.25\]
Substituting in equation \rfi{i},
\[A=40\,(1-e^{-2.25})=40\times(1-0.1054)=35.8\ \text{kg CO}_2\,\text{ha}^{-1}\,\text{h}^{-1}\Ans\]
\sh{2. Shaded leaf}
\[\frac{\varepsilon I}{A_{\max}}=\frac{0.45\times50}{40}=0.5625\]
\[A=40\,(1-e^{-0.5625})=40\times(1-0.5698)=17.2\ \text{kg CO}_2\,\text{ha}^{-1}\,\text{h}^{-1}\Ans\]
\sh{3. Efficiency of the use of light}
\[\frac AI\Bigr|_{\text{sunlit}}=\frac{35.8}{200}=0.179,\qquad\frac AI\Bigr|_{\text{shaded}}=\frac{17.2}{50}=0.344\]
\[\therefore\quad\frac{0.344}{0.179}=1.9\]
The shaded leaf uses its light nearly twice as efficiently, because it works on the linear part of the curve while the sunlit leaf works near saturation.

\section{Maintenance and growth respiration}\label{sec:dm-resp}
Part of the assimilate is burnt. \kw{Maintenance respiration} keeps the living tissue; it is proportional to the dry mass and doubles for each 10 $^\circ$C (a $Q_{10}$ of 2):
\begin{equation}\label{eq:dm-rm}\boxed{R_m=c_m\,W\,Q_{10}^{(T-T_{\mathrm{ref}})/10},\qquad \Delta W=E_g\,(A-R_m)}\end{equation}
The remainder is converted into new tissue with the \kw{conversion efficiency} $E_g$ (about 0.72 for leaves and stems, lower for oil-rich seeds); the loss $1-E_g$ is the \kw{growth respiration}. A large, old crop in a warm period can spend half of its daily assimilate on maintenance.

\exbox{The daily carbon budget of a crop}{A crop of 800 g\,m$^{-2}$ dry mass assimilates 25 g\,m$^{-2}$\,d$^{-1}$ of carbohydrate at a mean temperature of 15 $^\circ$C. Take $c_m=0.015$ g\,g$^{-1}$\,d$^{-1}$ at $T_{\mathrm{ref}}=25$ $^\circ$C, $Q_{10}=2$ and $E_g=0.72$. Find the maintenance respiration and the growth. Repeat at 25 $^\circ$C.}\sol
\givenline{$W=800$ g\,m$^{-2}$ ; $A=25$ g\,m$^{-2}$\,d$^{-1}$ ; $c_m=0.015$ g\,g$^{-1}$\,d$^{-1}$ at $T_{\mathrm{ref}}=25$ $^\circ$C ; $Q_{10}=2$ ; $E_g=0.72$}
\wkt\ (Equation~\ref{eq:dm-rm})
\[R_m=c_m\,W\,Q_{10}^{(T-T_{\mathrm{ref}})/10}\ti{i}\]
\[\Delta W=E_g\,(A-R_m)\ti{ii}\]
\sh{1. At 15 $^\circ$C}
From equation \rfi{i},
\[R_m=0.015\times800\times2^{(15-25)/10}=12.0\times2^{-1}=6.0\ \text{g\,m}^{-2}\,\text{d}^{-1}\Ans\]
and from equation \rfi{ii},
\[\Delta W=0.72\times(25-6.0)=0.72\times19.0=13.7\ \text{g\,m}^{-2}\,\text{d}^{-1}\Ans\]
\sh{2. At 25 $^\circ$C}
\[R_m=0.015\times800\times2^{0}=12.0\ \text{g\,m}^{-2}\,\text{d}^{-1}\Ans\]
\[\Delta W=0.72\times(25-12.0)=9.4\ \text{g\,m}^{-2}\,\text{d}^{-1}\Ans\]
\[\therefore\quad\text{Fall of growth}=\frac{13.7-9.4}{13.7}=0.31\]
A warm period reduces the growth by about a third through respiration alone, if the assimilation is unchanged; this is one reason why warming lowers the yield of temperate cereals.

\section{Water stress and yield: the yield-response factor}\label{sec:dm-ky}
Over a season the relative loss of yield is nearly proportional to the relative deficit of evapotranspiration. Doorenbos and Kassam (FAO paper 33) wrote
\begin{equation}\label{eq:dm-ky}\boxed{1-\frac{Y_a}{Y_m}=K_y\Bigl(1-\frac{ET_a}{ET_m}\Bigr)}\end{equation}
where $Y_a$ and $Y_m$ are the actual and the maximum yield and $K_y$ the \kw{yield-response factor}. $K_y>1$ (maize 1.25, wheat 1.15) means that the yield falls faster than the water use: such crops should not be stressed. $K_y<1$ (groundnut 0.7, soybean 0.85) allows a planned deficit. Daisy does not use Equation~\ref{eq:dm-ky}; its yield follows from the stress factor on photosynthesis day by day, but the twenty seasons of a SWAFL run can be plotted in this form.

\begin{kfig}{The yield-response factor: relative yield loss against relative evapotranspiration deficit. A crop with $K_y>1$ loses more yield than water.}
\draw[axis] (0,0)--(4.6,0);\draw[axis] (0,0)--(0,4.6);
\node[lab,below] at (2.2,-0.1) {$1-ET_a/ET_m$};\node[lab,rotate=90] at (-0.45,2.2) {$1-Y_a/Y_m$};
\draw[gd] (0,0)--(4.2,4.2);\node[lab] at (3.9,3.5) {$K_y=1$};
\draw[crv] (0,0)--(3.4,4.25);\node[lab,left] at (3.0,4.1) {$K_y=1.25$};
\draw[ln] (0,0)--(4.2,2.94);\node[lab] at (4.1,2.6) {$K_y=0.7$};
\end{kfig}

\exbox{Yield under a planned deficit}{A wheat crop ($K_y=1.15$) with a potential yield of 7.0 Mg\,ha$^{-1}$ receives irrigation that meets 80 per cent of its seasonal water use. Find the expected yield.}\sol
\givenline{$K_y=1.15$ ; $Y_m=7.0$ Mg\,ha$^{-1}$ ; $ET_a/ET_m=0.80$}
\wkt\ (Equation~\ref{eq:dm-ky})
\[1-\frac{Y_a}{Y_m}=K_y\Bigl(1-\frac{ET_a}{ET_m}\Bigr)\ti{i}\]
Substituting in equation \rfi{i},
\[1-\frac{Y_a}{Y_m}=1.15\times(1-0.80)=1.15\times0.20=0.23\]
\[\therefore\quad\frac{Y_a}{Y_m}=0.77\qquad\text{and}\qquad Y_a=0.77\times7.0=5.39\ \text{Mg\,ha}^{-1}\Ans\]
A 20 per cent deficit of water costs 23 per cent of the yield, because $K_y>1$.

\section{The distribution of the roots}\label{sec:dm-rootdist}
The root length density usually falls exponentially with depth, $L_v(z)=L_{v0}\,e^{-z/\lambda}$, with a decay length $\lambda$ of 15 to 30 cm for annual crops. The share of the roots above a depth $z_1$ in a root zone of depth $Z_r$ is
\begin{equation}\label{eq:dm-rootfrac}\boxed{F(z_1)=\frac{1-e^{-z_1/\lambda}}{1-e^{-Z_r/\lambda}}}\end{equation}
The uptake of Equation~\ref{eq:dm-feddes} is distributed in the same proportion when the soil is wet, which is why the upper soil dries first and why a sensor in the upper root zone sees the depletion early.

\exbox{Where are the roots?}{A crop has a root zone 100 cm deep with $\lambda=20$ cm. What fraction of its roots lies in the top 30 cm?}\sol
\givenline{$Z_r=100$ cm ; $\lambda=20$ cm ; $z_1=30$ cm}
\wkt\ the share of the roots above $z_1$ is (Equation~\ref{eq:dm-rootfrac})
\[F(z_1)=\frac{1-e^{-z_1/\lambda}}{1-e^{-Z_r/\lambda}}\ti{i}\]
Here $z_1/\lambda=30/20=1.5$ and $Z_r/\lambda=100/20=5$. Substituting in equation \rfi{i},
\[F(30)=\frac{1-e^{-1.5}}{1-e^{-5}}=\frac{1-0.2231}{1-0.0067}=\frac{0.7769}{0.9933}=0.782\Ans\]
Nearly four-fifths of the roots lie in the top 30 cm; a sensor at 20 cm therefore watches the part of the profile that delivers most of the water.

